\begin{answer}
We use the assumption $p(t=1\mid x)\in\{0,1\}$, so $h(x) = p(y=1|x) = \alpha\cdot p(t=1|x) \in \{0, \alpha\}$.

\textbf{Case $y^{(i)}=1$:} Observing $y=1$ requires $t=1$ (since $p(y=1|t=0,x)=0$). Hence $p(t=1|x^{(i)})=1$ and $h(x^{(i)}) = \alpha\cdot 1 = \alpha$.

\textbf{Case $y^{(i)}=0$:} From $p(y=1|x)=\alpha\cdot p(t=1|x)$, the observation $y=0$ is consistent with either $p(t=1|x)=0$ (giving $h(x)=0$) or $p(t=1|x)=1$ (giving $h(x)=\alpha$).

For the purpose of estimating $\alpha$, only the $y=1$ case is needed. Since $h(x^{(i)})=\alpha$ for every $x^{(i)}$ with $y^{(i)}=1$:
\[
\mathbb{E}[h(x)\mid y=1] = \alpha,
\]
which motivates the estimator $\hat\alpha = \frac{1}{|V_+|}\sum_{x^{(i)}\in V_+} h(x^{(i)})$. $\blacksquare$
\end{answer}
